\section{Long Run Performance}
\subsection{Classification of States}
\paragraph{Communicating Classes}
State $j$ is \textcolor{Bittersweet}{accessible from state $j$}, $i\to j$ if
\[\mathbf{P}_{ij}^{m}>0\text{ for some }m\ge 0.\]
If $i\to j\wedge j\to i$, they are \textcolor{Bittersweet}{communicating}.
$i\leftrightarrow j$.

Communication is an \textcolor{Blue}{equivalence relation}.
We define $(S/\leftrightarrow)$ the set of communication classes, each denoted
$\mathcal{C}$. It is \textcolor{Blue}{well-defined}.

Classes do not communicate. If $i\in \mathcal{C}_1,j\in \mathcal{C}_2$ such
that $i\to j$,
\[\forall i' \in \mathcal{C}_1, j' \in \mathcal{C}_2\qquad i'\to j',\quad
j\centernot\to i,\quad j'\centernot\to i'.\]

A MC is \textcolor{Bittersweet}{irreducible} if there is only one communication
class.

\paragraph{Recurrence}
For state $i$, the \textcolor{Bittersweet}{first return probability} is the
probability of having the $n$-th transition be the first return to $i$ with
initial state $i$
\[f_{ii}^{(0)}:=0.\] 
\[f_{ii}^{(n)}=P\left( X_1\neq i,\ldots,X_{n-1}\neq i, X_{n}=i|X_0=i \right).\] 
The probability of \textcolor{Bittersweet}{revisiting $i$} in the future is
\[f_{ii}=\sum_{n=0}^{\infty} f_{ii}^{(n)}\le 1,\]
and the \textcolor{Bittersweet}{return probability} to $i$ from $i$ at the
$n$-th transition is
\[\mathbf{P}_{ii}^{(n)}=P\left( X_{n}=i |X_0=i\right).\]
We have the \textcolor{Blue}{convolution}
\[P_{ii}^{(n)}=\sum_{k=1}^{n} f_{ii}^{(k)}P_{ii}^{(n-k)}.\]

The \textcolor{Bittersweet}{expected number of revisits} to state $i$, $E\left[
N_i|X_0=i \right]$ is
\begin{align*}
  E\left[ N_i|X_0=i \right]=&\sum_{n=1}^{\infty} P_{ii}^{(n)}\\
                            =&\begin{cases}\infty& f_{ii}=1\\\frac{f_{ii}}{1-f_{ii}}&f_{ii}<1\end{cases}
\end{align*}

\begin{align*}
  \text{$i$ is \textbf{\textcolor{Bittersweet}{recurrent/not recurrent}}} &\iff f_{ii}=1\\
                                                                   &\iff E\left[ N_i|X_0=i \right]\text{\ is infinite}\\
                                                                   &\iff \sum_{n=1}^{\infty} P_{ii}^{(n)}
\end{align*}
Recurrence is a \textcolor{Blue}{class property}.

\paragraph{Periodicity}
For a state $i$, we define the \textcolor{Bittersweet}{period} of $i$ 
\[d(i):=\begin{cases} \gcd{R_{ii}}&R_{ii}\neq \emptyset \\0&R_{ii}=\emptyset \end{cases}\]
where $R_{ij}:=\left\{ n\ge 1:P_{ij}^{(n)}>0 \right\}$.
\begin{itemize}
  \item $\exists N\in \mathbb{N}\quad\text{such that}\quad\forall n\ge N,0<k<d(i)$
    \[\mathbf{P}_{ii}^{(n*d(i))}>0,\qquad \mathbf{P}_{ii}^{(n*d(i)+k)}=0.\] 
  \item $\forall m\in R_{ji}$, for the same $N$, $\forall n\ge N$
    \[\mathbf{P}_{ji}^{(m+nd(i))}>0.\] 
\end{itemize}
Periodicity is a \textcolor{Blue}{class property}. A class with period 1 is
\textcolor{Bittersweet}{aperiodic}.

\paragraph{Regularity and Ergodicity}
A MC is \textcolor{Bittersweet}{regular} if $\exists k\in \mathbb{N}$ such that
\[\forall i,j\in S ,\quad {\left( \mathbf{P}^{(k)} \right)}_{ij}>0.\] 
$\mathbf{P}$ of a regular MC is a regular transition probability matrix.

\subsection{Recurrent Irreducible MC}
The \textcolor{Bittersweet}{expected duration between revisits} is defined 
\[m_i=E\left[ R_i|X_0=i \right] =\sum_{n=1}^{\infty} nf_{ii}^{(n)}.\] 
If $m_i$ is finite, the MC is \textcolor{Bittersweet}{positive recurrent}.
Otherwise, it is null recurrent. Given a fixed path, consider the
\textcolor{Bittersweet}{long-run proportion of visits}. It is given by $\lim_{n
  \to \infty} E\left[ \frac{1}{n}\sum_{k=1}^{n}
I\left( X_k=j|X_0=i \right) \right]$
\[\lim_{n \to \infty} \frac{1}{n}\sum_{k=1}^{n} \mathbf{P}_{ij}^{(k)}=\frac{1}{m_j}.\] 
If the MC is periodic:
$\lim_{n \to \infty} \mathbf{P}_{jj}^{(nd)}=\frac{d}{m_j}$.\\
If the MC is aperiodic and null recurrent:
$\lim_{n \to \infty} \mathbf{P}_{ij}^{(n)}=0$.\\
\textbf{\textcolor{ForestGreen}{Finite MC}}: Regular$\iff$ Ergodic $\iff$ Irreducible, Aperiodic.
\textbf{\textcolor{ForestGreen}{Infinite MC}}: Ergodic $\iff$ Irreducible,
Aperiodic, + Recurrent.

\paragraph{Ergodic MC}
$\lim_{n \to \infty} \mathbf{P}_{ij}^{(n)}$ exists independent of initial state
$i$.
\[\lim_{n \to \infty} \mathbf{P}_{ij}^{(n)}=\lim_{n \to \infty} \mathbf{P}_{jj}^{(n)}=\frac{1}{m_j}=\pi_j.\] 
Since $\sum_{k=1}^{N} \pi_k=1$, we define the \textcolor{Bittersweet}{limiting
distribution}
\[\mathbf{\pi} := \left<\pi_1,\pi_2,\ldots,\pi_N \right>.\] 
$\mathbf{\pi}$ is obtained from the solution to
$\mathbf{\pi}\mathbf{P}=\mathbf{\pi}$, $\left\lVert \mathbf{\pi} \right\rVert_1
=1$
\[\begin{cases}\pi_j=\sum_{k=1}^{N} \pi_k\mathbf{P}_{kj}&j\in \left\{ 1,2,\ldots,N \right\}\\
\left\lVert \pi \right\rVert_1 =1\end{cases}.\] 
The regular transition probability matrix converges
\[\lim_{n \to \infty} \mathbf{P}^{n}=\begin{pmatrix} \pi\\\pi\\\vdots\\\pi \end{pmatrix}=
\bordermatrix{&1&2&\cdots&N\cr1&\pi_1&\pi_2&\cdots&\pi_N\cr 2&\pi_1&\pi_2&\cdots&\pi_N\cr \vdots&\vdots&\vdots&\ddots&\vdots\cr N&\pi_1&\pi_2&\cdots&\pi_N}\] 
That is, $X_{n}\sim \pi^{(0)}\mathbf{P}^{(n)}\to\pi$ for any initial distribution.\\
Given a large time point $n$, \textcolor{Blue}{$X_{n}$ is a random variable}. Each path
gives a realisation of $X_{n}$ and $P(X_{n}=k)=\pi_j$.\\
Given a path, consider the \textcolor{Blue}{proportion of visits} over time. $E\left[ \frac{1}{n}\sum_{k=0}^{n-1} I\left( X_k=j|X_0=i \right) \right]\to\pi_j$ as $n\to \infty$.\\
The \textcolor{Blue}{mean duration between visits} to $j$ is given by $\frac{1}{m_j}$.

\subsection{General MC Long Run Performance}
\begin{enumerate}
  \item Obtain the set of communicating classes $\mathcal{C}_k$
  \item Set up a MC such that each recurrent class is denoted by a node
  \item Define $T:=\min{\left\{ n : X_{n}\in \text{recurrent classes}  \right\}} $
  \item Using FSA, find $u_{k|i} = P\left( X_T\in \mathcal{C}_k | X_0=i \right) $.
  \item Within each recurrent class $\mathcal{C}_k$:
    \begin{itemize}
      \item $\text{Period} > 1$: Long-run proportion of time in each state.
      \item $\text{Period}=1$: If finite or infinite positive recurrent,
        find the corresponding limiting distribution of state $j\in
        \mathcal{C}_k$, denoted $\pi_{j|k}$. Otherwise, the probability at each
        state is $0$.
    \end{itemize}
  \item For any initial state $X_0=i$:
    \begin{itemize}
      \item If $j$ is transient, $\pi_{j|i}=0$
      \item If $j\in \mathcal{C}_k$, $\pi_{j|i}=u_{k|i}\pi_{j|k}$ 
    \end{itemize}
  \item With an initial distribution $X_0\sim \pi^{(0)}$
    \[\pi_{j|\pi^{(0)}}=\sum_{i\in S} \pi_{j|i}\pi_{i}^{(0)}.\]
\end{enumerate}

\subsection{Balance Equations}
The \textcolor{Bittersweet}{global balance equations} $\pi=\pi\mathbf{P}$
yields the limiting distribution. Equivalently, we have the $N$ equations:
\[\pi_j=\sum_{k\in S} \pi_k\mathbf{P}_{kj}.\] 
$\pi$ exists even for non-ergodic MCs: $(X_{n+1}|X_{n}\sim\pi)\sim \pi$.
\vspace{-5mm}
\begin{figure}[H]
  \centering
  \includegraphics[width=0.17\textwidth]{stationary_dist}
\end{figure}
\vspace{-6mm}
A \textcolor{Blue}{unique} stationary distribution is the limiting distribution.\\
Consider a transition matrix $\mathbf{M}$ such that
\[\mathbf{M}_{ij}=P\left( X_{n+1}=j,X_{n}=i \right) =\pi_i\mathbf{P}_{ij},\]
The row sum equals the column sum for any $i,j\in S$.
This can be obtained with the $\binom{n}{2}$ \textcolor{Bittersweet}{local
(detailed) balance equations}. Suppose there exists a distribution $\pi$ such
that
\[\pi_i\mathbf{P}_{ij}=\pi_j\mathbf{P}_{ji}.\] 
Then matrix $\mathbf{M}$ is symmetric.
$\pi$ also \textcolor{Blue}{satisfies the global balance equations}, and is
easier to compute.
